\begin{theorem}[Physical range recovered from a regular cut]%
    \label{thm:peps_regular_ginjective_cut_range}
    \lean{TNLean.PEPS.regularBoundaryCutMatrix,
        TNLean.PEPS.IsGInjective.range_regularBoundaryCutMatrix,
        TNLean.PEPS.IsGInjective.rank_regularBoundaryCutMatrix,
        TNLean.PEPS.range_eq_of_regularBoundaryCutMatrix_eq}
    \leanok
    \uses{thm:peps_g_injective_invariants_range}
    Let $G$ be a finite group and let $T$ and $S$ be $G$-injective
    maps from the regular virtual boundary with $b\geq1$ legs to
    the two physical spaces of a cut. The coefficient matrix
    $C=TS^{\mathsf T}$ satisfies
    \[
        \operatorname{im}C=\operatorname{im}T,
        \qquad \operatorname{rank}C=|G|^{b-1}.
    \]
    Consequently, two such factorizations of the same coefficient
    matrix have the same first physical range, even when their
    virtual groups or numbers of boundary legs differ. This is a
    consequence of the regular boundary construction in
    \cite[Definition~5.1 and proof of Theorem~6.9]{Schuch2010PEPS}.
\end{theorem}

\begin{proof}\leanok
    \uses{thm:peps_g_injective_invariants_range,
        thm:peps_regular_boundary_invariants,
        thm:peps_regular_boundary_density}
    Write $P_G$ for the regular boundary averaging projector.
    There is a linear map $L$ with $LS=P_G$. Since
    $P_G^{\mathsf T}=P_G$
    (Theorem~\ref{thm:peps_regular_boundary_density}) and $TP_G=T$,
    one has $CL^{\mathsf T}=T$. The two contractions are
    % Formula: C_{p,q}=sum_x T_{p,x} S_{q,x};
    % sum_q C_{p,q} L_{x,q}=T_{p,x}.
    % Source: regular boundary cut in the proof of Theorem 6.9,
    % Papers/1001.3807/paper_v3.tex:1927--1984; this is the matrix
    % factorization used here, not a claim to reproduce the lattice geometry.
    % Ink-to-index: the T--S bond bundles all b crossing group labels into x.
    % C has physical indices p and q. L converts q back to virtual x.
    % Boundary signatures: the first row has (V,P)=(0,2), the second (1,1).
    % T and S share a virtual basis without conjugation: the factor is S^T.
    \begin{tenkzequation}
        \tenkzkernel
        \begin{tenkz}[rows={wire}, cols=1, bonds=none]
            \tn[skin=box, ports={180:physical, 0:physical}]{C}
        \end{tenkz}
        $=$
        \begin{tenkz}[rows={wire}, cols=2, bonds=none]
            \tn[name=cutLeft, ports={180:physical, 0:virtual}]{T}
            \tn[at=(1,2), name=cutRight,
                ports={180:virtual, 0:physical}]{S}
            \tnwire{cutLeft.0}{cutRight.180}
        \end{tenkz}
    \end{tenkzequation}
    \begin{tenkzequation}
        \tenkzkernel
        \begin{tenkz}[rows={wire}, cols=2, bonds=none]
            \tn[name=cutCoefficient, skin=box,
                ports={180:physical, 0:physical}]{C}
            \tn[at=(1,2), name=cutInverse, skin=box,
                ports={180:physical, 0:virtual}]{L}
            \tnwire{cutCoefficient.0}{cutInverse.180}
        \end{tenkz}
        $=$
        \begin{tenkz}[rows={wire}, cols=1, bonds=none]
            \tn[ports={180:physical, 0:virtual}]{T}
        \end{tenkz}
    \end{tenkzequation}
    Thus $\operatorname{im}T\subseteq\operatorname{im}C$,
    while $C=TS^{\mathsf T}$ gives the reverse
    inclusion. The rank follows from the dimension $|G|^{b-1}$ of the
    invariant boundary space
    (Theorem~\ref{thm:peps_regular_boundary_invariants}). Applying the
    range identity to each factorization proves the comparison assertion.
\end{proof}
